
Section: Many kinds of causality. The science of causality is the study of cause and effect relationships, and is a fundamental tool for understanding
and reasoning about the world and how it works. Correct causal reasoning is crucial to making correct decisions,
building robust systems, and scientific discovery itself. While advances in causal modeling have formalized core
concepts of causality, the varying tasks and problems across domains have led different fields to use related but distinct
conceptualizations and tools.
Here we describe three orthogonal categories of causal approaches and tasks. The first axis (covariance- and logic-based
causality) distinguishes methodological approaches that emphasize data analysis from those that emphasize logical
reasoning. The second axis (type vs. token causality) 2 focuses on the setting for a causal question: Are we asking
about causal relationships in general, such as for a population, or asking questions about specific causes for a specific
event. The third axis, presented in Section 2.2 focus on the causal question or task itself: are we interested in inferring
new relationships, characterizing the strength of specific known or hypothesized relationships, attributing blame or
reward for some outcome, etc. \textcite{kiciman2023causal}

